\section{结论与展望}
本文围绕复杂刚体在多体动力学中的接触建模问题，提出了基于符号距离场与压力场耦合的分布式顺应接触研究思路。本文认为，现有点接触模型虽然在计算上简洁高效，但在面接触、接触斑演化与分布式接触力矩表达方面存在先天局限；而完整连续体接触模型虽然物理上更充分，却往往不适于多体动力学中的轻量化计算。基于此，本文尝试寻找一条中间路径：利用 SDF 的复杂几何表达能力生成接触斑，并在接触斑上构造压力场，通过积分获得净接触力与力矩，从而建立一种兼具几何适应性、物理合理性和工程可实现性的接触建模方法。

围绕该方法，本文进一步给出了问题定义、接触斑构造准则、压力场闭合关系、数值实现流程以及基准验证协议。后续工作将重点围绕实验对比、参数辨识、自适应离散积分及摩擦扩展等方面展开，以形成完整的理论、算法与工程验证体系。
